% !TEX root = ../main.tex
\chapter{Analyse af eksisterende system}
\label{ch:analyse}

Før den praktiske integrationsindsats kan beskrives, er det nødvendigt at
forstå det system, der danner udgangspunkt for projektet. Dette kapitel
gennemgår den oprindelige Wattson-robot på systemniveau, beskriver den
arm, der udgør rapportens fokus, og redegør kort for den eksisterende
dokumentation og softwarestack. Til sidst begrundes valget af armen som
afgrænset delsystem.

\section{Robotten som helhed}

Wattson er en bipedal humanoid robot, hvor samtlige bærende mekaniske dele
er fremstillet ved hjælp af 3D-print, specifikt i PLA. Robotten består overordnet af følgende
delsystemer:

\begin{itemize}
    \item To arme, hver med syv frihedsgrader (skulder, overarm, albue,
          underarm og håndled)
    \item To hænder med adskilte finger-led
    \item Hoved med to frihedsgrader (pitch og yaw)
    \item Torso, hofte og to ben med flere led pr.\ ben
\end{itemize}

Robotten styres af en enkelt computer (Jetson Orin Nano), som
afvikler en eksisterende ROS~2 Humble-baseret softwarestack. Aktuatorerne
er fordelt på flere dedikerede mikrocontrollere (ESP32), der hver styrer
en logisk gruppe af servoer og kommunikerer med Jetsonen via USB-serial.
Den oprindelige robot understøtter teleoperation via et VR-headset og et
ZED-stereo\-kamera, men disse funktioner er ikke en del af rapportens scope
(jf.\ kapitel~\ref{ch:indledning}).

% TODO: indsæt et samlet block-diagram her -- robotten med alle delsystemer
% og deres forbindelser til Jetsonen via ESP32-controllere.
\begin{figure}[H]
    \centering
    \begin{tikzpicture}[
        scale=0.72, transform shape,
        node distance=0.9cm and 0.2cm,
        every node/.style={font=\small},
        host/.style={draw, rectangle, rounded corners=2pt,
            minimum height=0.8cm, minimum width=4.4cm,
            align=center, fill=blue!8, thick},
        block/.style={draw, rectangle, rounded corners=2pt,
            minimum height=0.8cm, minimum width=3.6cm,
            align=center, fill=gray!8},
        sensor/.style={draw, rectangle, rounded corners=2pt,
            minimum height=0.8cm, minimum width=3.2cm,
            align=center, fill=green!8},
        servos/.style={draw, rectangle, dashed,
            minimum height=0.9cm, minimum width=3.6cm,
            align=center, fill=gray!3},
        line/.style={-{Latex[length=1.6mm]}, thick},
    ]
        \node[host] (jetson) {Jetson Orin Nano\\\footnotesize ROS\,2 Humble \textbar{} 4$\times$ USB};

        \node[sensor, above=0.7cm of jetson] (zed) {ZED-stereokamera};

        \node[block, below left=1.0cm and 0.1cm of jetson] (esp1)
            {ESP32 -- gul};
        \node[block, below=1.0cm of jetson] (esp2)
            {ESP32 -- rød};
        \node[block, below right=1.0cm and 0.1cm of jetson] (esp3)
            {ESP32 -- hvid};

        \node[servos, below=of esp1] (s1)
            {Venstre arm\\\footnotesize 7$\times$ ST3215};
        \node[servos, below=of esp2] (s2)
            {Højre arm + hoved\\\footnotesize 7$\times$ ST3215 + 2$\times$ servo};
        \node[servos, below=of esp3] (s3)
            {Hænder\\\footnotesize 10$\times$ SC09};

        \draw[line] (zed) -- node[right, font=\scriptsize]{USB} (jetson);
        \draw[line] (jetson.south) -- (esp1.north);
        \draw[line] (jetson.south) -- (esp2.north);
        \draw[line] (jetson.south) -- (esp3.north);
        \draw[line] (esp1) -- (s1);
        \draw[line] (esp2) -- (s2);
        \draw[line] (esp3) -- (s3);
    \end{tikzpicture}
    \caption{Overordnet styringsarkitektur for Wattson-robotten.
    Farverne svarer til den fysiske kabelmærkning,
    jf.~kapitel~\ref{ch:el}.}
    \label{fig:robot-overview}
\end{figure}

\section{Armen}
\label{sec:arm-delsystem}

En arm er den enhed, der i rapporten behandles som selvstændigt delsystem.
Armen består af syv aktive led, der tilsammen udgør en antropomorf
kinematik: skulderen har tre frihedsgrader (pitch, roll og yaw), albuen
én (pitch), underarmen én (yaw), og håndleddet to (pitch og roll).

\subsection{Aktuatorer}

De syv led i armen drives af Waveshare ST3215 seriebus-servoer.
Dette er en intelligent servotype med indbygget mikrocontroller, magnetisk
positionsencoder og en intern PID-løkke. Servoerne kommunikerer over en
halv-duplex TTL-seriel bus, herefter blot servobussen, hvor flere
servoer er forbundet i daisy-chain og adresseres ved hvert sit
unikke ID.

Den interne PID-løkke gør, at den overordnede styring kun behøver sende
en ønsket position til hver servo; reguleringen mod målpositionen sker
lokalt i servoen. Reguleringsmekanismen er beskrevet i
kapitel~\ref{ch:el}, og selve PID-parametrene indlæses fra en
JSON-konfigurationsfil, der er dokumenteret i
afsnit~\ref{sec:doks-kode}.

\subsection{Styreelektronik}

Servoerne er ikke forbundet direkte til Jetsonen, men til en dedikeret
ESP32-mikrocontroller, der fungerer som bus-master på den serielle
servobus. ESP32'en er via USB forbundet til Jetsonen og videresender
kommandoer fra Jetsonen ud til de enkelte servoer.

I den fuldt opbyggede robot er der tre ESP32-controllere, én pr.\
logisk gruppe. De er mærket med farverne gul, rød og hvid, der svarer
til kabelfarverne (jf.\ kapitel~\ref{ch:el}). I den afgrænsede
konfiguration anvendes kun den, der styrer den valgte arm.

% TODO: kort beskrivelse af ESP32'ens firmware -- protokol mod Jetsonen.
% Skal undersøges/afklares med Nicoline. Hvis tiden ikke rækker, kan det
% beskrives på højt niveau: "ESP32'en oversætter beskeder fra ROS-noden
% til ST3215-protokollen og videresender feedback den anden vej."

\clearpage
\subsection{Mekanisk opbygning}

De mekaniske led består af 3D-printede dele, der monteres direkte
på servoernes aksler. Belastningen i leddet bæres altså af både
selve den printede del og af servoens lejer. Denne konstruktion er let og
billig, men introducerer nogle særlige hensyn omkring styrke, stivhed og
slid, som behandles i kapitel~\ref{ch:diskussion}.

\subsection{Signal- og effektkæde}

På et overordnet niveau ser signal- og effektkæden for én arm ud som
vist i Figur~\ref{fig:arm-blokdiagram}. Signal- og effektvejene er
adskilte og mødes først ved den enkelte servo: positionskommandoer kommer
fra Jetsonen via USB og servobussen, mens forsyningen kommer direkte fra
strømforsyningen via en WAGO-klemrække. Den nominelle driftsspænding
% TODO: bekræft eksakt spænding mod den anvendte PSU; kapitel~\ref{ch:el}
% angiver 12\,V via Mean~Well~DDR-120.
for ST3215 er ca.\ 12\,V; den eksakte forsynings\-værdi fremgår af
kapitel~\ref{ch:el}.

\begin{figure}[H]
    \centering
    \begin{tikzpicture}[
        node distance=0.8cm and 0.9cm,
        every node/.style={font=\small},
        host/.style={draw, rectangle, rounded corners=2pt,
            minimum height=0.9cm, minimum width=2.8cm,
            align=center, fill=blue!8, thick},
        block/.style={draw, rectangle, rounded corners=2pt,
            minimum height=0.9cm, minimum width=2.4cm,
            align=center, fill=gray!8},
        psu/.style={draw, rectangle, rounded corners=2pt,
            minimum height=0.9cm, minimum width=2.8cm,
            align=center, fill=orange!12, thick},
        servo/.style={draw, rectangle, rounded corners=2pt,
            minimum height=1.1cm, minimum width=2.4cm,
            align=center, fill=gray!3, thick},
        sigline/.style={-{Latex[length=2mm]}, thick},
        pwrline/.style={-{Latex[length=2mm]}, thick, double, double distance=1pt},
    ]
        % Signal chain (top row)
        \node[host] (jetson) {Jetson Orin Nano};
        \node[block, right=of jetson] (esp) {ESP32\\\footnotesize bus-master};

        % Power chain (bottom row)
        \node[psu, below=1.6cm of jetson] (psu) {Strømforsyning};
        \node[block, right=of psu] (rail) {WAGO-klemrække};

        % Servo group (right)
        \node[servo, right=1.4cm of esp,
              anchor=west, yshift=-0.9cm,
              minimum height=2.6cm] (servos)
            {7$\times$ ST3215\\\footnotesize ID 1--7\\(daisy-chain)};

        % Signal arrows
        \draw[sigline] (jetson) -- node[above, font=\scriptsize]{USB} (esp);
        \draw[sigline] (esp.east) -- node[above, font=\scriptsize]{servobus}
            (servos.west |- esp.east);

        % Power arrows
        \draw[pwrline] (psu) -- (rail);
        \draw[pwrline] (rail.east) -- node[above, font=\scriptsize]{$\sim$12\,V}
            (servos.west |- rail.east);

        % Legend
        \node[anchor=north west, font=\scriptsize, align=left]
            at ([yshift=-0.4cm]psu.south west)
            {\tikz\draw[sigline] (0,0) -- (0.6,0); \ signal \quad
             \tikz\draw[pwrline] (0,0) -- (0.6,0); \ effekt};
    \end{tikzpicture}
    \caption{Signal- og effektkæde for én arm. Signal- og effektveje er
    adskilte og mødes først ved servoerne. Spændingsangivelsen er
    nominel; den faktiske forsynings\-værdi for den valgte PSU fremgår
    af kapitel~\ref{ch:el}.}
    \label{fig:arm-blokdiagram}
\end{figure}
\clearpage
\section{Eksisterende dokumentation og kildekode}
\label{sec:doks-kode}

Hele Wattson-projektet er offentligt tilgængeligt som et ROS~2-workspace
på GitHub. Workspacet indeholder bl.a.\ pakker til servo-styring,
teleoperation, kamera-integration og opstart af robotten. For rapportens
formål er især følgende relevant:

\begin{itemize}
    \item \enquote{servos.launch.py} -- starter en
          \enquote{wattson\_servo\_manager\_node}, der indlæser
          servokonfigurationer fra en angiven mappe og opretter
          ROS-topics for hver gruppe.
    \item \enquote{slider\_control.launch.py} -- en GUI med skyderkontroller
          pr.\ led. Den gør det enkelt at sende styresignaler direkte
          til den enkelte servo og dermed verificere, at hvert led
          reagerer som forventet.
    \item \enquote{servo\_arm\_left\_params.json} -- konfigurationsfil for
          den venstre arms syv servoer. Indeholder for hvert led bl.a.\
          servo-ID, gear-ratio, PID-gains (\enquote{gain\_P},
          \enquote{gain\_I}, \enquote{gain\_D}), software-grænser og
          udgangsposition.
\end{itemize}

Det er væsentligt, at konfigurationsfilen for en arm er selvstændig.
Hver arm kan derfor i princippet køres uafhængigt af resten af robotten,
hvilket er en forudsætning for det modulære arbejde i denne rapport.
Den enkelte fils opbygning og de ændringer, der er foretaget under
projektet, er dokumenteret i Bilag~\ref{AppendixC},
afsnit~\ref{sec:appB-json}.

\section{Begrundelse for valg af arm som delsystem}

Armen er valgt som afgrænset delsystem, fordi den i kompakt form
indeholder alle de elementer, der karakteriserer Wattson-projektet som
helhed:

\begin{itemize}
    \item Et 3D-printet skelet med flere bevægelige led
    \item Flere elektriske aktuatorer, der skal forsynes og styres
          samtidigt.
    \item En dedikeret mikrocontroller (ESP32), der skal idriftsættes
          og kommunikere med en overordnet computer.
    \item En reguleringsmæssig styring -- både den lokale PID-løkke i
          hver servo og den overordnede positionskommando fra ROS.
    \item En kommunikationskæde i flere niveauer (USB, servobus,
          servo-protokol).
\end{itemize}

\noindent Armen rummer dermed både den mekaniske, elektriske og
styringsmæssige side af projektet i én samlet enhed.
